\clearpage
\phantomsection

\addcontentsline{toc}{chapter}{{KẾT LUẬN VÀ HƯỚNG NGHIÊN CỨU TIẾP THEO}}

\chapter*{Kết luận và hướng nghiên cứu tiếp theo}

\section*{1. Tổng kết kết quả}

\begin{table}[H]
    \centering
    \small
    \caption{Kết quả đạt được theo từng mục tiêu đề tài}
    \label{tab:mt_results_summary}
    \begin{tabular}{|c|p{4.5cm}|p{4.5cm}|p{2.5cm}|}
        \hline
        \textbf{MT} & \textbf{Mục tiêu} & \textbf{Kết quả} & \textbf{Run tốt nhất} \\
        \hline
        MT1 & Agent chiến đấu hiệu quả trong Roguelike & Clear rate 0.626, kills 1.53/episode & run62 \\
        \hline
        MT2 & Giải quyết phần thưởng thưa + bản đồ thay đổi & Curriculum +85\% so với mốc so sánh (2.073→3.843) & run46, run50 \\
        \hline
        MT3 & Xác định kỹ thuật hỗ trợ tác động nhất & LSTM/ICM/RND kém hơn curriculum đơn thuần & run46 vs run35/36/42 \\
        \hline
        MT4 & Kiểm tra nút thắt quan sát (observation bottleneck) & Cấu hình run62 phá plateau 4.40→3.10 nat & run62 \\
        \hline
    \end{tabular}
\end{table}

Khóa luận đã xây dựng và đánh giá một tác nhân học tăng cường sâu cho môi trường game hành động Roguelike trong Unity. Hệ thống cuối cùng sử dụng CombatAgent với quan sát vector 207 chiều, không gian hành động rời rạc nhiều nhánh $(3,3,3,9)$, pipeline sinh bản đồ bằng PCGNN và cơ chế curriculum learning điều khiển qua tham số difficulty. Bản đồ được sinh ngoại tuyến, chuyển sang định dạng văn bản của game, kiểm tra khả năng đi tới bằng BFS và chia thành Easy, Medium, Hard để phục vụ curriculum.

Kết quả thực nghiệm cho thấy curriculum learning là xu hướng nổi bật nhất trong nhóm ablation chính về kỹ thuật huấn luyện. PPO baseline trên phân phối khó đạt phần thưởng tail-50 2.073, trong khi run46 với curriculum đạt trung bình ở Hard 3.843, tương đương mức cải thiện khoảng 85\%. Khi ghép ICM lên nền curriculum, run50 đạt trung bình ở Hard 4.369 và 1.016 kills/episode, gợi ý rằng phần thưởng nội tại có thể cộng hưởng khi tác nhân đã được hỗ trợ bằng các lesson dễ hơn. Biến thể fixed-to-random spawn trong run54 đạt trung bình ở Hard 3.925, cho thấy việc tổ chức độ khó môi trường là một yếu tố nền quan trọng. Ở nhóm ablation mở rộng về biểu diễn quan sát, run59 với đặc trưng \emph{path-informed} tạo ra bước cải thiện thực nghiệm lớn nhất; bản cập nhật run61 đạt tỷ lệ dọn phòng cuối run khoảng 0.846, số kẻ địch tiêu diệt cuối run khoảng 1.78 và entropy giảm khỏi vùng plateau quanh 4.40 nat xuống khoảng 3.62 ở vùng cuối.

Cuối cùng, run62 kết hợp quan sát lưới cục bộ, curriculum 4 bậc, ICM và \texttt{aggressive\_seek}, sau đó được huấn luyện đến 33.45M bước. Run này đạt phần thưởng tail-50 14.105, tỷ lệ dọn phòng tail-50 0.626, số kẻ địch tiêu diệt tail-50 1.530 và entropy tail-50 3.099, phá vỡ hoàn toàn entropy plateau 4.40 nat quan sát ở các ablation 1--3M bước với quan sát 207 chiều. Kết quả này gợi ý mạnh rằng plateau không phải giới hạn cứng của PPO hay action space, nhưng vì run62 thay đổi nhiều yếu tố đồng thời, nguyên nhân cụ thể chưa được cô lập hoàn toàn. Các kết quả toàn bộ chưa đi kèm error bar do mỗi cấu hình chỉ được chạy một seed, vì vậy cần được diễn giải như bằng chứng thực nghiệm ban đầu thay vì kết luận thống kê cuối cùng.

\begin{table}[H]
    \centering
    \small
    \caption{Các kết quả nổi bật của khóa luận}
    \label{tab:conclusion_best_results}
    \begin{tabular}{|p{4.5cm}|c|p{7cm}|}
        \hline
        \textbf{Kết quả} & \textbf{Giá trị} & \textbf{Ý nghĩa} \\
        \hline
        PPO baseline & 2.073 phần thưởng tail-50 & Mốc so sánh không curriculum. \\
        \hline
        Run46 (curriculum) & 3.843 trung bình ở Hard & Tăng khoảng 85\% so với mốc so sánh trong run đã thực hiện. \\
        \hline
        Run50 (curriculum + ICM) & 4.369 trung bình ở Hard & Cải thiện thêm trên nền curriculum, nhưng inverse loss vẫn cao. \\
        \hline
        Run54 (fixed-to-random) & 3.925 trung bình ở Hard & Gần run46, cho thấy fixed-to-random spawn là một tinh chỉnh môi trường có tác động nhỏ hơn curriculum. \\
        \hline
        Run55 (reward v2) & 0.800 đỉnh tỷ lệ dọn phòng & Tăng giá trị đỉnh của tỷ lệ dọn phòng, nhưng phần thưởng trung bình không so trực tiếp do dùng thang thưởng khác. \\
        \hline
        Run59/run61 (\emph{path-informed}) & tỷ lệ dọn phòng cuối run khoảng 0.846, số kẻ địch tiêu diệt cuối run khoảng 1.78 & Bước đột phá về biểu diễn quan sát; bổ sung tín hiệu tìm đường giúp tác nhân tiếp cận và tiêu diệt kẻ địch ổn định hơn. \\
        \hline
        Run62 (lưới cục bộ, 33.45M bước) & 14.105 phần thưởng tail-50, 0.626 dọn phòng, entropy 3.099 & Phá vỡ entropy plateau; tỷ lệ dọn phòng tiệm cận BFS mà không cài cứng cơ chế tìm đường. \\
        \hline
        Entropy plateau & 4.40 nat (ablation 1--3M bước); 3.099 nat (run62, 33.45M bước) & Plateau bị phá vỡ trong run62, khi kết hợp quan sát lưới cục bộ, curriculum, ICM và 33.45M bước huấn luyện. \\
        \hline
    \end{tabular}
\end{table}

Các thí nghiệm ICM, RND và LSTM đem lại thông tin bổ sung nhưng không làm thay đổi xu hướng chính. Khi chạy độc lập, ICM có dấu hiệu inverse model học yếu trên action space nhiều nhánh; RND cải thiện nhẹ hơn ICM nhưng không vượt trần; LSTM có xu hướng học chậm và tăng dần nhưng chưa tạo thay đổi nổi bật trong quy mô huấn luyện hiện tại. Khi ghép với curriculum, LSTM không tạo cộng hưởng rõ, còn ICM trong run50 giúp tăng phần thưởng và số lần tiêu diệt nhưng inverse loss vẫn cao.

\section*{2. Đóng góp của khóa luận}

Đóng góp thứ nhất là xây dựng được một pipeline thực nghiệm hoàn chỉnh cho RL trong game hành động Roguelike: từ thiết kế quan sát/hành động, sinh bản đồ bằng PCGNN, chuẩn hóa bản đồ sang text asset, kiểm tra khả năng đi tới bằng BFS, đến huấn luyện PPO bằng Unity ML-Agents. Pipeline này không chỉ chạy trên một bản đồ cố định mà mở rộng sang tập bản đồ được chia theo độ khó.

Đóng góp thứ hai là kết quả thực nghiệm về curriculum learning trong môi trường chiến đấu sinh theo thủ tục. Trong phạm vi một seed, đưa tác nhân qua chuỗi Easy -- Medium -- Hard giúp chính sách học được kỹ năng nền trước khi đối mặt với phân phối khó, tạo cải thiện lớn hơn so với phần thưởng nội tại hoặc thay đổi kiến trúc bộ nhớ khi các kỹ thuật này chạy độc lập.

Đóng góp thứ ba là pipeline tuyển chọn bản đồ PCGNN cho curriculum. Hệ thống đã bổ sung bộ chuyển đổi đặt \texttt{P}/\texttt{E}, kiểm tra đường đi, tính các chỉ số như tỷ lệ tường, độ dài đường đi, tỷ lệ ngõ cụt, A* difficulty và chia bản đồ theo percentile. Với lô 14x14 gồm 1000 bản đồ, hệ thống tạo được 1000 bản đồ duy nhất, 100\% bản đồ có đường đi hợp lệ và chia thành 50 Easy, 50 Medium, 900 Hard để người thiết kế lựa chọn.

Đóng góp thứ tư là các kết quả ablation có giá trị. ICM khi chạy độc lập không hoạt động tốt trong action space multi-discrete của bài toán, thể hiện qua inverse loss quanh 4.36 so với mức ngẫu nhiên lý thuyết 5.49. Khi ghép với curriculum, run50 cải thiện phần thưởng Hard lên 4.369 nhưng inverse loss vẫn ở 4.337, gợi ý rằng curiosity có thể hỗ trợ hành vi nhưng chưa giải quyết triệt để vấn đề biểu diễn.

Đóng góp thứ năm là phát hiện và giải quyết entropy plateau. Trong phạm vi các ablation chính 207 chiều, entropy duy trì quanh 4.40 nat bất kể thay đổi ICM, RND, LSTM, curriculum, giảm phương sai vị trí sinh hoặc thiết kế lại phần thưởng. Cấu hình run62 sau đó cho thấy plateau này bị phá vỡ khi kết hợp quan sát lưới cục bộ, curriculum 4 bậc, ICM và 33.45M bước huấn luyện: entropy giảm xuống 3.099 nat và duy trì bền vững, cho thấy vấn đề nằm ở biểu diễn quan sát và quy mô huấn luyện, không phải ở kiến trúc PPO hay action space.

Đóng góp thứ sáu là ablation mở rộng về biểu diễn quan sát thông qua run59. Khi bổ sung đặc trưng \emph{path-informed} dựa trên BFS, tác nhân đạt bước cải thiện lớn nhất về hành vi: tỷ lệ dọn phòng, số kẻ địch tiêu diệt và mức độ quyết đoán của chính sách đều tốt hơn rõ rệt. Kết quả này cho thấy biểu diễn không gian là yếu tố quyết định trong môi trường chiến đấu sinh theo thủ tục, và mở ra hướng kết hợp RL với các đặc trưng điều hướng có cấu trúc thay vì chỉ thay đổi thuật toán huấn luyện.

\section*{3. Hạn chế}

Hạn chế lớn nhất của khóa luận là mỗi run mới chỉ chạy một seed. Vì vậy, các kết luận về chênh lệch nhỏ cần được hiểu là xu hướng thực nghiệm, chưa phải kết quả thống kê có error bar. Các phát hiện có chênh lệch lớn, như curriculum tăng khoảng 85\%, đáng chú ý hơn so với các chênh lệch nhỏ; tuy nhiên vẫn cần chạy nhiều seed để xác nhận độ ổn định.

Hạn chế thứ hai là PCGNN hiện chạy ngoại tuyến và có bước con người tuyển chọn. Cách này giúp ổn định huấn luyện và giảm lỗi tích hợp, nhưng chưa phải một cơ chế sinh nội dung thích ứng theo thời gian thực, nơi độ khó bản đồ thay đổi trực tiếp dựa trên năng lực hiện tại của tác nhân. Ngoài ra, khóa luận chưa có một đánh giá tách biệt đủ sâu về khả năng khái quát hóa trên bản đồ chưa từng thấy, chẳng hạn so sánh tác nhân huấn luyện trên bản đồ thủ công với tác nhân huấn luyện bằng curriculum PCGNN rồi kiểm thử trên tập bản đồ PCGNN chưa từng thấy.

Hạn chế thứ ba là quy mô huấn luyện giữa các nhóm chưa hoàn toàn đồng nhất. Nhóm không curriculum chủ yếu dừng ở 1.0M bước, trong khi nhóm curriculum chạy đến 3.0M bước và run62 đạt 33.45M bước. Vì vậy, một số so sánh phản ánh hiệu quả của cấu hình hoàn chỉnh hơn là ablation tuyệt đối công bằng theo cùng số bước.

Một hạn chế bổ sung được rút ra từ run59 là bài toán điều hướng trong bản đồ sinh theo thủ tục chưa được giải quyết hoàn toàn bằng quan sát cục bộ 207 chiều. Khi bổ sung hướng bước tiếp theo theo BFS và khoảng cách đường đi ngắn nhất vào tín hiệu học, hiệu năng tăng rõ. Vì vậy, run59 được xem là một ablation mở rộng về biểu diễn quan sát, cho thấy yếu tố tìm đường có tác động lớn tới khả năng tiếp cận và tiêu diệt kẻ địch, đồng thời gợi ý rằng điều hướng là một nút thắt quan trọng của môi trường hiện tại.

Hạn chế cuối cùng là nguyên nhân chính xác dẫn đến entropy plateau chưa được cô lập. Run62 đã phá vỡ plateau bằng cách kết hợp nhiều thay đổi đồng thời (lưới cục bộ, curriculum 4 bậc, ICM, hyperparameter điều chỉnh và số bước huấn luyện lớn hơn), nhưng chưa có ablation đơn biến nào xác định yếu tố nào thực sự là nguyên nhân chính. Ngoài ra, các hướng như giảm số hướng aim, gộp nhánh combat--aim, thay đổi hệ số entropy, factorized policy hoặc khởi tạo bằng imitation learning chưa được kiểm tra.

\section*{4. Hướng nghiên cứu tiếp theo}

Hướng đầu tiên là thiết kế lại action space. Không gian $(3,3,3,9)$ giúp biểu diễn điều khiển game đầy đủ, nhưng cũng làm entropy cao và khiến chính sách khó trở nên quyết đoán. Một hướng khả thi là gộp nhánh combat/aim, giảm số hướng aim, hoặc dùng action mask mạnh hơn theo ngữ cảnh.

Hướng thứ hai là chạy lại các run quan trọng với nhiều seed. Tối thiểu nên chạy run34, run46, run50, run54 và run55 với ba seed để có trung bình, độ lệch chuẩn và khoảng tin cậy. Điều này sẽ giúp kết luận về curriculum, ICM trên nền curriculum và reward v2 có sức thuyết phục thống kê cao hơn.

Hướng thứ ba là tiếp tục khai thác quan sát lưới cục bộ. Run62 cho thấy khi số bước huấn luyện tăng lên 33.45M, lưới cục bộ có thể đạt tỷ lệ dọn phòng tiệm cận BFS-assisted mà không cần cài cứng cơ chế tìm đường. Hướng tiếp theo là tăng độ phân giải lưới, kết hợp với quan sát pixel hoặc đặc trưng đồ thị để kiểm tra khả năng khái quát hóa sang các bản đồ chưa từng thấy.

Hướng thứ tư là dùng imitation learning hoặc behavior cloning để khởi tạo chính sách ban đầu. Nếu có một tập minh họa nhỏ từ người chơi, tác nhân có thể bắt đầu từ chính sách biết di chuyển, né và tấn công cơ bản, sau đó dùng PPO/curriculum để tinh chỉnh trên phân phối bản đồ khó.

Hướng thứ năm là đưa PCGNN từ chế độ ngoại tuyến sang curriculum thích ứng. Khi tác nhân đạt hiệu năng cao ở một nhóm, hệ thống có thể sinh thêm bản đồ mới quanh vùng độ khó cao hơn, sau đó kiểm tra bằng chỉ số định lượng và đưa vào tập huấn luyện. Đây là hướng gần với sinh nội dung thủ tục thích ứng (adaptive PCG) và có thể giúp giảm hiện tượng tác nhân học quá khớp vào tập bản đồ cố định.

Hướng thứ sáu là tăng quy mô huấn luyện có hệ thống. Run62 đạt 33.45M bước và vẫn cho thấy xu hướng cải thiện; chưa rõ điểm bão hòa nằm ở đâu. Huấn luyện thêm kết hợp với đánh giá định kỳ trên tập bản đồ giữ riêng (held-out) sẽ giúp xác định khi nào nên dừng hoặc chuyển sang tinh chỉnh kiến trúc.

Hướng thứ bảy là mở rộng sang huấn luyện multi-agent. Môi trường TopDown Engine hỗ trợ nhiều nhân vật và nhiều kẻ địch, vì vậy các kịch bản hợp tác hoặc đối kháng sẽ là hướng tự nhiên để kiểm tra khả năng phối hợp, chia mục tiêu và thích nghi chiến thuật của tác nhân.

Mã nguồn và tài liệu liên quan được cập nhật tại:
\url{https://github.com/tranmanh2004}
